\documentclass[12pt,a4paper]{article}
\usepackage[margin=1in]{geometry}
\usepackage{graphicx}
\usepackage{listings}
\usepackage{booktabs}
\usepackage{array}
\usepackage{hyperref}
\usepackage{xcolor}
\lstset{basicstyle=\ttfamily\small,breaklines=true,frame=single,columns=fullflexible}
\hypersetup{colorlinks=true,urlcolor=blue,linkcolor=black}
\setlength{\parskip}{0.6em}

\title{\textbf{Repository / Directory Synchronizer}\\[4pt]
\large A Client--Server Application for Synchronizing a Local Folder with a Remote Folder}
\author{CSE 302 -- Operating Systems, Section N\\
Group 7, Project ID 4\\[4pt]
V.~Preetham (AP24110010605), K.~Aadarsh (AP24110010702),\\
A.~Akhil (AP24110010633), M.~Jathin (AP24110010131),\\
M.~Sasikanth (AP24110010551), T.~Jaideep (AP24110011909)\\[4pt]
SRM University-AP}
\date{\today}

\begin{document}
\maketitle

\begin{abstract}
Keeping the same set of files on two machines is a common need, and tools such as Dropbox and \texttt{rsync} solve it. This project implements a small version of such a tool: a client--server application that watches a local directory and synchronizes every change to a remote directory. The client detects file creation, modification, and deletion, hashes file contents with SHA-256 to avoid sending unchanged data, and sends changes to the server over TCP. The server uses a thread pool, and a mutex protects its shared index. The project uses core operating system concepts: threads, mutual exclusion, the producer--consumer pattern, inter-process communication through sockets, and file system metadata. The source code is available on GitHub and a web-based simulation of the algorithm is hosted on GitHub Pages.
\end{abstract}

\section{Introduction}
A user who works on files on one computer often needs a copy of them on another. Copying files by hand is slow and error-prone, and copying everything every time wastes bandwidth. A synchronizer solves this by detecting what changed and transferring only that.

The goal of this project is to build a client--server synchronizer that copies local changes to a remote folder, as the Dropbox desktop client does. The project also serves as a practical exercise in operating system concepts that appear in the CSE 302 syllabus.

\subsection{Problem Statement}
Design and implement a client--server application that synchronizes changes made in a local folder to a remote folder, and that handles several simultaneous requests safely.

\subsection{Objectives}
\begin{itemize}
\item Detect file creation, modification, and deletion in a local directory automatically.
\item Transfer only files whose content has changed.
\item Serve multiple connections at the same time without corrupting shared state.
\item Recover from network failures by retrying.
\item Reject unsafe requests such as path traversal.
\end{itemize}

\section{Related Work}
\textbf{Dropbox} keeps a local folder and a cloud folder consistent, uploads changes in the background, and keeps copies of conflicting versions.\
\textbf{rsync} compares files between two locations and transfers only the differences, using rolling checksums to find the changed parts of a file.\
\textbf{Git} identifies content by hash, which is the same idea this project uses to decide whether a file has changed.

Our project is much smaller than these systems. It keeps the central ideas (change detection, hashing, background transfer) and leaves out features such as delta encoding and user accounts.

\section{System Architecture}
The system has two programs. The client runs next to the folder being watched, and the server stores the remote copy.

\begin{figure}[h]
\centering
\begin{lstlisting}
 CLIENT                                      SERVER
 +-----------+  event  +---------+   TCP    +--------------+
 | Watcher   | ------> |  Queue  | -------> | Thread pool  |
 | (producer)|         | (blocks)|          | (8 workers)  |
 +-----------+         +----+----+          +------+-------+
                            | consumer             | mutex
                       +----v-----+         +------v-------+
                       | Uploader |         | Index + files|
                       | SHA-256  |         | remote_store |
                       +----------+         +--------------+
\end{lstlisting}
\caption{Architecture of the synchronizer.}
\end{figure}

\subsection{Client}
The client has three parts:
\begin{enumerate}
\item A \textbf{watcher thread} (the \texttt{watchdog} library) that receives file system events and puts them on a queue.
\item A \textbf{blocking queue} that connects the watcher to the uploader.
\item An \textbf{uploader thread} that takes events from the queue, computes a SHA-256 hash, and sends the file to the server if the hash differs from the last synchronized one.
\end{enumerate}
When the client starts, it also walks the whole directory and queues every file, so that files created while the client was off are synchronized.

\subsection{Server}
The server accepts TCP connections on port 9000 and passes each connection to a worker in a thread pool of eight threads. A worker reads a one-line JSON header (operation, path, size), then the file contents, and updates the remote folder and the in-memory index.

\subsection{Protocol}
Each request starts with one line of JSON, for example \texttt{\{"op":"put","path":"a/b.txt","size":120\}}, followed by \texttt{size} bytes of file data. The operations are \texttt{put} and \texttt{delete}. The server replies with \texttt{OK}, or \texttt{ERR} if the request was rejected.

\section{Operating System Concepts Used}
\begin{table}[h]
\centering
\begin{tabular}{@{}p{3.6cm}p{11cm}@{}}
\toprule
\textbf{Concept} & \textbf{How it is used}\\
\midrule
Threads & The client runs a watcher thread and an uploader thread. The server runs a pool of worker threads, so one slow transfer does not block the others.\\
Mutual exclusion & Several workers update the shared index and write files at once. A mutex makes each update atomic and prevents a race condition.\\
Producer--consumer & The watcher produces events and the uploader consumes them. The queue blocks when empty, so the uploader does not busy-wait and wastes no CPU.\\
Inter-process communication & Client and server are separate processes that communicate through TCP sockets.\\
File system & The client traverses directories and reads files. The server validates paths and creates directories as needed.\\
\bottomrule
\end{tabular}
\caption{Mapping of OS concepts to the implementation.}
\end{table}

\section{Algorithms}
\subsection{Change Detection by Hashing}
A file's modification time can change without its content changing, for example when an editor saves the same text. The client therefore compares SHA-256 hashes.

\begin{lstlisting}
function process(event):
    if event is delete:
        remove path from synced; send DELETE
    else:
        h = SHA256(read(path))
        if synced[path] == h: return      // unchanged
        synced[path] = h; send PUT(path, data)
\end{lstlisting}

\subsection{Producer--Consumer Queue}
\begin{lstlisting}
watcher:   on event e:  queue.put(e)
uploader:  loop: e = queue.get()   // blocks while empty
                 process(e)
\end{lstlisting}

\subsection{Retry with Exponential Backoff}
If a send fails, the uploader waits $2^{k}$ seconds before attempt $k+1$, for up to five attempts. This handles a server that is temporarily down without flooding it with requests.

\subsection{Conflict Resolution}
A conflict occurs when the same file is changed on both sides between two synchronizations. The web simulation demonstrates the policy: both versions are kept, and the second one is saved under the name \texttt{name (conflict).ext}. No data is lost, and the user decides which version to keep.

\subsection{Path Validation}
The server normalizes every received path and rejects paths that start with \texttt{..} or are absolute, which prevents a client from writing outside the remote folder.

\section{Implementation}
The project is written in Python 3. The server uses the \texttt{socket}, \texttt{threading}, and \texttt{concurrent.futures} modules. The client uses \texttt{watchdog} for file system events, \texttt{queue.Queue} for the blocking queue, and \texttt{hashlib} for SHA-256.

\begin{table}[h]
\centering
\begin{tabular}{@{}ll@{}}
\toprule
\textbf{File} & \textbf{Purpose}\\
\midrule
\texttt{src/server.py} & Threaded TCP server, index, storage\\
\texttt{src/client.py} & Watcher, queue, uploader, retry\\
\texttt{docs/index.html} & Web simulation hosted on GitHub Pages\\
\texttt{report/main.tex} & This report\\
\bottomrule
\end{tabular}
\caption{Repository contents.}
\end{table}

To run the system: start the server with \texttt{python src/server.py}, then start the client with \texttt{python src/client.py ./my\_folder}.

\section{Testing and Results}
Each test was run with the server and client on the same machine. The \emph{Result} column records what was observed (Pass or Fail).

\begin{table}[h]
\centering
\begin{tabular}{@{}p{0.4cm}p{4.2cm}p{6.4cm}p{1.6cm}@{}}
\toprule
\textbf{\#} & \textbf{Test} & \textbf{Expected behaviour} & \textbf{Result}\\
\midrule
1 & Create a new file & File appears in \texttt{remote\_store} & \_\_\_\\
2 & Modify a file & Remote file has the new content & \_\_\_\\
3 & Delete a file & File is removed from \texttt{remote\_store} & \_\_\_\\
4 & Save a file without changes & No upload happens (hash matches) & \_\_\_\\
5 & Stop the server, change a file, restart the server & Client retries and the change arrives & \_\_\_\\
6 & Request path \texttt{../x} & Server replies \texttt{ERR}, nothing is written & \_\_\_\\
\bottomrule
\end{tabular}
\caption{Test cases.}
\end{table}

% Add screenshots of the two terminals and the two folders here:
% \begin{figure}[h]\centering
% \includegraphics[width=0.9\textwidth]{screenshot1.png}
% \caption{Client and server output during a sync.}\end{figure}

\section{Limitations and Future Work}
\begin{itemize}
\item Synchronization is one-way (local to remote). A pull from remote to local would make it bidirectional.
\item Whole files are sent. Splitting files into chunks and sending only changed chunks would save bandwidth on large files.
\item The index is in memory, so it is lost when the server restarts.
\item There is no encryption or authentication.
\item File renames are seen as a delete followed by a create.
\end{itemize}

\section{Conclusion}
The project delivers a working client--server synchronizer that detects local changes, avoids sending unchanged files, and serves requests with a thread pool protected by a mutex. It shows how threads, mutual exclusion, the producer--consumer pattern, and sockets work together in a real application. The source code and the web simulation are linked below.

\section*{Links}
GitHub repository: \url{https://github.com/AP2411001190/dir-synchronizer}\\
Web app: \url{https://ap2411001190.github.io/dir-synchronizer/}

\begin{thebibliography}{9}
\bibitem{osc} A.~Silberschatz, P.~B.~Galvin, G.~Gagne, \emph{Operating System Concepts}, Wiley.
\bibitem{rsync} A.~Tridgell, P.~Mackerras, ``The rsync algorithm,'' Australian National University, 1996.
\bibitem{dropbox} Dropbox Developers, \url{https://www.dropbox.com/developers}.
\bibitem{watchdog} \texttt{watchdog} Python library, \url{https://pypi.org/project/watchdog/}.
\bibitem{py} Python Software Foundation, Python documentation: \texttt{socket}, \texttt{threading}, \texttt{hashlib}, \texttt{queue}.
\end{thebibliography}

\end{document}
